\documentclass[11pt,a4paper]{article}
\usepackage[utf8]{inputenc}
\usepackage[T1]{fontenc}
\usepackage[french]{babel}
\usepackage[a4paper,top=1.75cm,bottom=1.9cm,left=1.85cm,right=1.85cm]{geometry}
\usepackage{amsmath,amssymb}
\usepackage{array}
\usepackage{booktabs}
\usepackage{enumitem}
\usepackage{eurosym}
\usepackage{fancyhdr}
\usepackage{hyperref}
\usepackage{lmodern}
\usepackage{inconsolata}
\usepackage{listings}
\usepackage{microtype}
\usepackage{tabularx}
\usepackage[most]{tcolorbox}
\usepackage{titlesec}
\usepackage{xcolor}

\renewcommand{\familydefault}{\sfdefault}

\definecolor{ink}{HTML}{202124}
\definecolor{muted}{HTML}{5F6368}
\definecolor{line}{HTML}{DADCE0}
\definecolor{paper}{HTML}{F8F9FA}
\definecolor{panel}{HTML}{FFFFFF}
\definecolor{accent}{HTML}{8A1538}
\definecolor{accentsoft}{HTML}{F8EEF2}
\definecolor{green}{HTML}{0B6B57}
\definecolor{greensoft}{HTML}{E6F4EF}
\definecolor{amber}{HTML}{9A5B00}
\definecolor{ambersoft}{HTML}{FFF7E6}
\definecolor{blue}{HTML}{245B8E}
\definecolor{bluesoft}{HTML}{EAF2FA}
\definecolor{rose}{HTML}{A84D77}
\definecolor{rosesoft}{HTML}{FAF0F5}
\definecolor{codeback}{HTML}{F8F9FA}
\definecolor{codemark}{HTML}{DADCE0}

\hypersetup{
  colorlinks=true,
  linkcolor=accent,
  urlcolor=green,
  pdfborder={0 0 0}
}

\setlength{\parindent}{0pt}
\setlength{\parskip}{5pt}
\setlist[itemize]{leftmargin=1.35em,itemsep=2pt,topsep=2pt}
\setlist[enumerate]{leftmargin=1.55em,itemsep=3pt,topsep=3pt}
\setlist[enumerate,1]{label=\textbf{\alph*.}}

\pagestyle{fancy}
\fancyhf{}
\fancyhead[L]{\sffamily\footnotesize Python pour économistes}
\fancyhead[R]{\sffamily\footnotesize Université de Lille -- L2 Économie}
\fancyfoot[C]{\sffamily\footnotesize\color{muted}\thepage}
\renewcommand{\headrulewidth}{0.2pt}
\renewcommand{\headrule}{\hbox to\headwidth{\color{line}\leaders\hrule height \headrulewidth\hfill}}
\setlength{\headheight}{22pt}

\titleformat{\section}
  {\sffamily\Large\bfseries\color{ink}}
  {}{0pt}{}
  [\vspace{-0.25em}{\color{accent}\titlerule[0.7pt]}]
\titleformat{\subsection}
  {\sffamily\large\bfseries\color{ink}}
  {}{0pt}{}

\lstdefinelanguage{terminal}{
  morekeywords={ls,cd,pwd,mkdir,python,python3,code,dir,type,echo},
  sensitive=true
}

\lstset{
  basicstyle=\ttfamily\footnotesize,
  backgroundcolor=\color{codeback},
  frame=single,
  framesep=5pt,
  framerule=0.25pt,
  rulecolor=\color{codemark},
  language=Python,
  keywordstyle=\color{accent}\bfseries,
  commentstyle=\color{muted},
  stringstyle=\color{green},
  numberstyle=\scriptsize\color{muted},
  numbers=left,
  numbersep=8pt,
  showstringspaces=false,
  breaklines=true,
  aboveskip=0.7em,
  belowskip=0.7em,
  xleftmargin=1.1em,
  framexleftmargin=1em,
  columns=fullflexible,
  keepspaces=true,
  upquote=true,
  literate=
   {à}{{\`a}}1 {â}{{\^a}}1 {ä}{{\"a}}1
   {ç}{{\c c}}1
   {é}{{\'e}}1 {è}{{\`e}}1 {ê}{{\^e}}1 {ë}{{\"e}}1
   {î}{{\^\i}}1 {ï}{{\"\i}}1
   {ô}{{\^o}}1 {ö}{{\"o}}1
   {ù}{{\`u}}1 {û}{{\^u}}1 {ü}{{\"u}}1
   {À}{{\`A}}1 {Â}{{\^A}}1 {Ä}{{\"A}}1
   {Ç}{{\c C}}1
   {É}{{\'E}}1 {È}{{\`E}}1 {Ê}{{\^E}}1 {Ë}{{\"E}}1
   {Î}{{\^I}}1 {Ï}{{\"I}}1
   {Ô}{{\^O}}1 {Ö}{{\"O}}1
   {Ù}{{\`U}}1 {Û}{{\^U}}1 {Ü}{{\"U}}1
   {œ}{{\oe}}1 {Œ}{{\OE}}1
   {–}{{--}}1 {—}{{---}}1
}

\newcommand{\code}[1]{\texttt{#1}}
\newcommand{\pp}{\,points de pourcentage}
\newcommand{\taglabel}[1]{\scriptsize\sffamily\bfseries #1}

\newcounter{exercise}
\newcounter{extension}
\newcounter{logic}

\newcommand{\workbooktitle}[4]{%
  \setcounter{exercise}{0}%
  \setcounter{extension}{0}%
  \setcounter{logic}{0}%
  \fancyhead[L]{\sffamily\footnotesize #1 -- #2}%
  \begingroup\raggedright
  {\sffamily\Large\bfseries\color{ink}#1 -- #2\par}
  \vspace{0.25em}
  {\sffamily\small\color{muted}Python pour économistes -- Université de Lille, L2 Économie\par}
  \vspace{0.85em}
  {\sffamily\bfseries\color{accent}Question fil rouge.} #3\par
  \vspace{0.35em}
  {\sffamily\bfseries\color{green}Livrables.} #4\par
  \vspace{0.9em}
  {\color{line}\hrule height 0.45pt}
  \vspace{0.85em}
  \endgroup
}

\newenvironment{objectivebox}[1]{%
  \subsection*{#1}
  \vspace{-0.55em}
}{\vspace{0.15em}}

\newenvironment{routebox}[1]{%
  \subsection*{#1}
  \vspace{-0.45em}
}{\vspace{0.15em}}

\newenvironment{deliverablebox}[1]{%
  \subsection*{#1}
  \vspace{-0.45em}
}{\vspace{0.15em}}

\newenvironment{methodbox}[1]{%
  \subsection*{#1}
  \vspace{-0.45em}
}{\vspace{0.15em}}

\newenvironment{pathwaybox}[1]{%
  \par\addvspace{0.8em}
  {\sffamily\bfseries\color{blue}#1\par}
  \vspace{-0.25em}
  \begingroup\leftskip=0.85em\relax
  \noindent\color{ink}
}{\par\endgroup\addvspace{0.25em}}

\newenvironment{pseudocodeexercise}[1]{%
  \refstepcounter{logic}%
  \begin{tcolorbox}[
    enhanced,breakable,
    title={Logique \thelogic{} -- #1 \hfill\taglabel{PSEUDO-CODE}},
    colback=bluesoft,
    colframe=blue!55!black,
    coltitle=ink,
    colbacktitle=bluesoft,
    fonttitle=\sffamily\bfseries,
    boxrule=0.35pt,
    arc=0pt,
    left=10pt,right=10pt,top=7pt,bottom=7pt,
    before skip=8pt,after skip=8pt,
    borderline west={3pt}{0pt}{blue}
  ]%
}{\end{tcolorbox}}

\newenvironment{mandatoryexercise}[1]{%
  \refstepcounter{exercise}%
  \par\addvspace{1.05em}
  {\color{line}\hrule height 0.35pt}
  \vspace{0.55em}
  {\sffamily\large\bfseries\color{ink}Exercice \theexercise{} -- #1\par}
  {\taglabel{OBLIGATOIRE}\par}
  \vspace{0.25em}
}{\par\addvspace{0.65em}}

\newenvironment{extensionexercise}[1]{%
  \refstepcounter{extension}%
  \begin{tcolorbox}[
    enhanced,breakable,
    title={Pour aller plus loin \theextension{} -- #1},
    colback=ambersoft,
    colframe=amber!65!black,
    coltitle=ink,
    colbacktitle=ambersoft,
    fonttitle=\sffamily\bfseries,
    boxrule=0.35pt,
    arc=0pt,
    left=10pt,right=10pt,top=8pt,bottom=8pt,
    before skip=8pt,after skip=8pt,
    borderline west={3pt}{0pt}{amber}
  ]%
}{\end{tcolorbox}}

\newenvironment{checkpointbox}[1]{%
  \subsection*{#1}
  \vspace{-0.45em}
}{\vspace{0.15em}}

\newtcolorbox{takeawaybox}[1]{
  enhanced,breakable,
  title={#1},
  colback=greensoft,
  colframe=green,
  coltitle=ink,
  colbacktitle=greensoft,
  fonttitle=\sffamily\bfseries,
  boxrule=0.35pt,
  arc=0pt,
  left=10pt,right=10pt,top=8pt,bottom=8pt,
  before skip=10pt,after skip=8pt,
  borderline west={3pt}{0pt}{green}
}


\begin{document}

\workbooktitle
  {TD4}
  {Graphiques et simulation simple à partir d'OWID}
  {Comment utiliser Python pour visualiser une distribution et comprendre intuitivement la stabilisation d'une moyenne ?}
  {Un fichier \code{td4\_nom\_prenom.py}, deux figures \code{.png} et une interprétation courte.}

\begin{objectivebox}{Objectifs du TD}
\begin{itemize}
  \item Construire un graphique simple avec \code{matplotlib}.
  \item Utiliser \code{random.choice} pour simuler des tirages.
  \item Observer comment une moyenne simulée évolue quand le nombre de tirages augmente.
  \item Identifier une \code{SyntaxError} et une erreur d'indentation dans une boucle.
\end{itemize}
\end{objectivebox}

\begin{checkpointbox}{Progression du TD}
Les exercices 1 à 11 sont obligatoires. La simulation est volontairement simple : elle sert à apprendre les boucles, \code{random.choice} et les graphiques, pas à faire de statistique avancée.
\end{checkpointbox}

\begin{checkpointbox}{Source des données}
Les taux de fécondité sont des extraits pédagogiques arrondis d'Our World in Data.
\begin{lstlisting}[language=terminal]
https://ourworldindata.org/grapher/children-per-woman-un.csv
\end{lstlisting}
\end{checkpointbox}

\begin{pathwaybox}{Bibliothèques utilisées}
\code{matplotlib} sert à tracer les graphiques. \code{random} sert à tirer des valeurs au hasard. Si un import échoue, le problème vient probablement de l'environnement Python ou de la bibliothèque installée, pas de la logique de votre script.
\end{pathwaybox}

\begin{pseudocodeexercise}{Visualiser puis simuler}
\begin{lstlisting}
# Objectif : observer une distribution puis simuler des tirages
# Entree : liste de taux de fecondite
# 1. Calculer la moyenne observee
# 2. Tracer un graphique en barres
# 3. Tirer au hasard une valeur dans la liste
# 4. Repeter le tirage plusieurs fois
# 5. Calculer la moyenne des tirages
# 6. Observer si la moyenne se stabilise quand n augmente
\end{lstlisting}
\end{pseudocodeexercise}

\section*{A. Préparer les données et un graphique}

\begin{mandatoryexercise}{Créer l'extrait OWID}
Créer un dictionnaire où les clés sont des pays ou régions et les valeurs des taux de fécondité arrondis.
\begin{lstlisting}
fertilite = {
    "World": 2.3,
    "France": 1.8,
    "United States": 1.6,
    "Japan": 1.2,
    "India": 2.0,
    "Nigeria": 4.5,
    "Niger": 6.1,
    "Brazil": 1.6,
    "Ethiopia": 3.9,
    "China": 1.0
}
\end{lstlisting}
Créer ensuite la liste des valeurs et calculer la moyenne initiale. Ce nom sera réutilisé dans la simulation.
\begin{lstlisting}
valeurs = list(fertilite.values())
moyenne_initiale = sum(valeurs) / len(valeurs)
\end{lstlisting}
Afficher le nombre d'observations et la moyenne simple de ces valeurs.
\end{mandatoryexercise}

\begin{mandatoryexercise}{Observer clés et valeurs}
Avant le graphique, vérifier ce que renvoient les conversions suivantes.
\begin{lstlisting}
pays = list(fertilite.keys())
valeurs = list(fertilite.values())

print(pays[:3])
print(valeurs[:3])
print(len(pays) == len(valeurs))
\end{lstlisting}
Écrire en commentaire pourquoi le graphique en barres a besoin de deux listes de même longueur.
\end{mandatoryexercise}

\begin{mandatoryexercise}{Graphique en barres}
Avec \code{matplotlib}, tracer un graphique en barres des taux de fécondité.
\begin{lstlisting}
import matplotlib.pyplot as plt

pays = list(fertilite.keys())
valeurs = list(fertilite.values())

plt.figure()
plt.bar(pays, valeurs)
plt.xticks(rotation=45, ha="right")
plt.ylabel("Enfants par femme")
plt.title("Taux de fecondite - extrait OWID")
plt.tight_layout()
plt.savefig("fertilite_barres.png", dpi=300)
\end{lstlisting}
\end{mandatoryexercise}

\begin{mandatoryexercise}{Modifier un paramètre graphique}
Reprendre le graphique en barres et faire deux modifications seulement :
\begin{enumerate}
  \item changer le titre pour indiquer qu'il s'agit d'un extrait pédagogique ;
  \item sauvegarder la figure sous le nom \code{fertilite\_barres\_version2.png}.
\end{enumerate}
Vérifier que l'ancien fichier et le nouveau fichier existent tous les deux. Si un seul fichier existe, vérifier le nom passé à \code{plt.savefig(...)}.
\end{mandatoryexercise}

\begin{mandatoryexercise}{Interpréter le graphique}
Ajouter trois commentaires :
\begin{enumerate}
  \item quel pays a le taux le plus élevé dans l'extrait ?
  \item quel pays a le taux le plus faible ?
  \item pourquoi ce graphique est-il plus lisible qu'une liste brute ?
\end{enumerate}
\end{mandatoryexercise}

\section*{B. Simulation très guidée}

\begin{mandatoryexercise}{Un tirage aléatoire}
Importer \code{random}, puis tirer une valeur au hasard dans la liste des taux.
\begin{lstlisting}
import random

tirage = random.choice(valeurs)
print(tirage)
\end{lstlisting}
Exécuter plusieurs fois le script : le résultat est-il toujours le même ?
\end{mandatoryexercise}

\begin{mandatoryexercise}{Prédire le rôle de \code{range}}
Avant de simuler 20 tirages, prédire combien de lignes seront affichées par chaque boucle.
\begin{lstlisting}
for i in range(3):
    print(i)

for i in range(1):
    print("tirage")
\end{lstlisting}
Exécuter ensuite les deux fragments. Ajouter une phrase qui explique pourquoi \code{range(20)} permet de répéter exactement 20 tirages.
\end{mandatoryexercise}

\begin{mandatoryexercise}{Moyenne de plusieurs tirages}
Compléter le code suivant pour calculer la moyenne de 20 tirages.
\begin{lstlisting}
tirages = []
for i in range(20):
    tirages.append(random.choice(valeurs))

moyenne_tirages = sum(tirages) / len(tirages)
print(moyenne_tirages)
\end{lstlisting}
Comparer cette moyenne à la moyenne de la liste initiale.
\textbf{Mini-débogage.}
Avant de continuer, expliquer pourquoi \code{for i in range(20)} sans deux-points produit une \code{SyntaxError}, et pourquoi une ligne non indentée après \code{for} produit une erreur d'indentation.
\end{mandatoryexercise}

\begin{mandatoryexercise}{Stabilisation de la moyenne}
Calculer la moyenne des tirages pour plusieurs tailles : \code{5}, \code{20}, \code{100}, \code{1000}. Afficher les résultats dans une phrase.
\begin{lstlisting}
tailles = [5, 20, 100, 1000]
for n in tailles:
    tirages = []
    for i in range(n):
        tirages.append(random.choice(valeurs))
    print(n, sum(tirages) / len(tirages))
\end{lstlisting}
\end{mandatoryexercise}

\section*{C. Visualiser la simulation}

\begin{pseudocodeexercise}{Passer des tirages à une courbe}
\begin{lstlisting}
# Objectif : voir si la moyenne simulee se stabilise
# Entree : plusieurs tailles de simulation
# 1. Creer une liste vide pour les moyennes
# 2. Pour chaque taille n :
#       faire n tirages
#       calculer la moyenne des tirages
#       ajouter cette moyenne a la liste
# 3. Tracer tailles en abscisse et moyennes en ordonnee
# Sortie : graphique de stabilisation
\end{lstlisting}
\end{pseudocodeexercise}

\begin{mandatoryexercise}{Graphique de stabilisation}
Créer deux listes : \code{tailles} et \code{moyennes\_simulees}. Compléter d'abord le squelette, puis tracer le graphique.
\begin{lstlisting}
tailles = [5, 20, 100, 1000]
moyennes_simulees = []

for n in tailles:
    tirages = []
    for i in range(n):
        tirages.append(random.choice(valeurs))
    moyenne_n = ...
    moyennes_simulees.append(moyenne_n)

plt.figure()
plt.plot(tailles, moyennes_simulees, marker="o")
plt.axhline(moyenne_initiale, linestyle="--")
plt.xlabel("Nombre de tirages")
plt.ylabel("Moyenne simulée")
plt.tight_layout()
plt.savefig("simulation_moyenne.png", dpi=300)
\end{lstlisting}
La variable \code{moyenne\_initiale} doit avoir été calculée avant le graphique.
\end{mandatoryexercise}

\begin{mandatoryexercise}{Interprétation économique}
Ajouter quatre commentaires :
\begin{enumerate}
  \item que représente la moyenne de l'extrait ?
  \item que montre la simulation quand le nombre de tirages augmente ?
  \item pourquoi ce résultat ne décrit-il pas directement la population mondiale ?
  \item pourquoi la sélection des pays est-elle importante ?
\end{enumerate}
\end{mandatoryexercise}

\section*{D. Entraînement facultatif}

\begin{extensionexercise}{Pour aller plus loin, sans obligation}
Choisir au moins deux questions si le parcours obligatoire est terminé.
\begin{enumerate}
  \item Fixer la graine avec \code{random.seed(123)} et observer l'effet.
  \item Faire 10 simulations de taille 20 et comparer les moyennes obtenues.
  \item Tracer un histogramme des 1000 tirages.
  \item Supprimer \code{Niger} de la liste et recalculer la moyenne.
  \item Ajouter une ligne pour la médiane de la liste initiale.
  \item Écrire une fonction \code{moyenne\_tirages(n)}.
  \item Tester des tailles \code{10}, \code{50}, \code{500}, \code{5000}.
  \item Calculer l'écart entre moyenne simulée et moyenne initiale.
  \item Lire intuitivement la loi des grands nombres : que veut dire “stabiliser” ?
  \item Extension avancée : construire un intervalle \code{moyenne +/- 2 * ecart\_type / racine(n)} avec une formule donnée.
\end{enumerate}
\end{extensionexercise}

\begin{takeawaybox}{À retenir}
\begin{itemize}
  \item Un graphique rend une comparaison plus lisible qu'une liste brute de valeurs.
  \item Un import donne accès à une bibliothèque ou à un module non disponible par défaut.
  \item Une simulation répète une opération contrôlée pour observer un comportement.
  \item Dans une boucle, les deux-points et l'indentation font partie de la logique du programme.
\end{itemize}
\end{takeawaybox}

\end{document}
